\begin{center}
{\large\bfseries Préparation au Concours National Commun --- Informatique MP}\\[8pt]
\rule{0.9\linewidth}{0.5pt}\\[6pt]
{\Large\bfseries Épreuve d'Informatique}\\[4pt]
{\bfseries Filière MP --- Durée : 4 heures}\\[6pt]
{\itshape Voyage d'un paquet : de l'adresse IP au cadenas du navigateur}\\[4pt]
{\itshape Les protocoles TCP/IP et HTTP/HTTPS}\\[6pt]
\rule{0.9\linewidth}{0.5pt}\\[8pt]
{\small Auteur~: \href{https://orcid.org/0000-0002-6108-7176}{\textbf{Pr Agrégé El Hadiq Zouhair}}}\\[3pt]
{\small \href{https://cpgeacademy.org}{\textbf{cpgeacademy.org}} \quad---\quad Tous droits réservés \textcopyright\ cpgeacademy.org}
\end{center}
\medskip
\begin{center}
\fcolorbox{gray}{gray!5}{\begin{minipage}{0.93\linewidth}\small
\textbf{Avis important.} Ce document est une \textbf{initiative personnelle} du professeur,
à but strictement pédagogique. Il ne s'agit \textbf{pas} d'un sujet officiel du Concours
National Commun~; il s'en inspire uniquement par la forme et l'esprit.
\smallskip

\textbf{Sources.} Le sujet s'appuie sur les documents normatifs de l'\textsc{ietf} et sur
deux articles fondateurs~: \rfc{791} (IP), \rfc{4632} (CIDR), \rfc{1071} et \rfc{1624}
(somme de contrôle), \rfc{6298} (temporisateur de retransmission), \rfc{5681} (contrôle de
congestion), \rfc{9110} et \rfc{9112} (HTTP/1.1), \rfc{9111} (cache), \rfc{8446}
(TLS~1.3)~; V.~Jacobson, \emph{Congestion Avoidance and Control} (SIGCOMM~1988) et
M.~Mathis \emph{et al.}, \emph{The Macroscopic Behavior of the TCP Congestion Avoidance
Algorithm} (CCR~1997). \textbf{Aucune connaissance préalable en réseaux n'est requise}~:
tout ce qui est nécessaire est défini dans le sujet.
\smallskip

\textbf{Pour aller plus loin.} Les élèves peuvent traiter une \textbf{nouvelle version} de
ce problème dotée d'un \textbf{autocorrecteur des réponses}, et suivre un \textbf{cours
complet}, sur \href{https://cpgeacademy.org}{\textbf{cpgeacademy.org}}.
\end{minipage}}
\end{center}
\medskip
